\section{导读与能力目标}

\subsection{手册怎么用}
本手册与 GitHub 仓库配套使用：
\begin{itemize}
  \item \textbf{自学入口}：仓库根目录 \texttt{README.md}（周计划、Overleaf 同步说明）。
  \item \textbf{正文}：本 LaTeX 手册（\texttt{latex/}），可在 Overleaf 在线改、编译 PDF。
  \item \textbf{模板与清单}：\texttt{templates/}、\texttt{assets/reading-list/}。
\end{itemize}

\subsection{入门后应具备的能力}
完成 4--8 周入门路径后，期望你能够：
\begin{enumerate}
  \item 用关键词在 Google Scholar / 顶会官网检索文献，并维护一份结构化文献 list；
  \item 按「泛读→筛选→精读→笔记」读完领域综述与 1--2 篇经典论文；
  \item 用 PyTorch 跑通官方入门示例，并读懂一份规范开源仓库的目录结构；
  \item 独立完成一次组会汇报（论文汇报或周汇报），并保留实验/会议记录；
  \item 说清自己课题的问题定义、输入输出、评价指标与数据来源（哪怕仍很初步）。
\end{enumerate}

\subsection{建议的时间节奏（概览）}
\begin{center}
\begin{tabular}{clp{8cm}}
\toprule
阶段 & 时长 & 重点 \\
\midrule
环境与框架 & 第 1 周 & PyTorch 60 分钟教程、开发环境、Claude Code 入门 \\
文献方法 & 第 2--3 周 & 检索渠道、综述精读、文献 list、Zotero \\
经典模型 & 第 4--5 周 & ResNet / Transformer 等精读 + 代码浏览 \\
课题对齐 & 第 6--7 周 & 组内课题范式、数据集摸底、第一次正式汇报 \\
巩固输出 & 第 8 周 & 复现小实验、整理笔记、制定下阶段 roadmap \\
\bottomrule
\end{tabular}
\end{center}

更细的逐周 checklist 见仓库 \texttt{README.md}；可复制 \path{templates/learning_plan.md} 填写个人计划。

\subsection{三个一定（贯穿全程）}
\begin{enumerate}
  \item \textbf{论文读完一定要做笔记}：方法归纳、贡献、未解决问题、对自己课题的启发。
  \item \textbf{实验跑完一定要有记录}：时间、软硬件、代码版本、数据、结果、分析、下一步。
  \item \textbf{组会一定要有纪要 + 待办}：把老师和同学指出的问题落到下一阶段任务。
\end{enumerate}
